\chapter{Project management}
\label{chap:project-management}
The project followed an Agile-inspired development approach, where work was completed in short iterations with frequent review and adaptation. Rather than following a rigid waterfall model, the team prioritised flexibility, allowing requirements and priorities to evolve as the project progressed. This approach was suitable because the team was developing a novel system and requirements became clearer during implementation.

\section{Project Planning}
At the beginning of the project, the team held an initial planning session to discuss potential project ideas and define the overall goals and scope. Several different system concepts were evaluated based on feasibility, project requirements and the available development timeframe. After discussion, the team conducted a vote to select the most suitable project idea. The final decision was to develop a ticketing system.

Following the selection of the project topic, the team began outlining the structure and functionality of the system. Diagrams and flowcharts were created to visualise the system process and understand how different components of the system would interact. These diagrams assisted in planning the development process and clarifying the workflow. During this stage, the team also identified the Minimum Viable Product (MVP), which included the essential features necessary for the system to operate effectively. Defining the MVP helped prioritise development tasks and ensured that core functionality could be completed within the project deadline.

A project timeline was then established to divide the work into several phases, including planning, development, testing and documentation. This structured approach allowed the team to manage the workload more effectively and maintain steady progress throughout the project.

\section{Team Agreement and Working Practices}
At the start of the project, the team established a formal agreement to define roles, responsibilities, communication methods and expectations. Decisions were made through discussion and majority vote, with polling used if consensus was not possible. Each member was responsible for writing readable, maintainable code and ensuring their commits did not disrupt the repository. The team resolved merge conflicts collaboratively and conducted code reviews together.

The team established clear communication guidelines and used scheduled meetings, WhatsApp and Trello to coordinate work, track tasks and monitor progress. Trello was used as a Kanban-style task management board, where tasks moved through stages such as “To Do”, “In Progress”, and “Completed”. This approach allowed team members to track progress in real time. 

Procedures were also established to address missed deadlines, work that did not meet expected standards and temporary unavailability due to illness or other commitments. Conflicts were handled internally within the team whenever possible, with lecturers or teaching assistants involved only if necessary. The agreement emphasised respect, accountability and maintaining a positive team environment, while prioritising quality over quantity in feature development.

\section{Team Roles and Responsibilities}
Responsibilities within the team were distributed based on individual strengths and interests. Hafsa acted as the team leader and chaired all scheduled meetings, ensuring that the project stayed on track. Amey managed the repository and DevOps tasks, including maintaining the version control system and overseeing code integration. Jasmeen took responsibility for recording meeting minutes, documenting decisions and maintaining records of progress. All team members collectively managed the Trello board, updating tasks, monitoring deadlines and coordinating work throughout the project.

Although specific roles were assigned, the project remained highly collaborative. Team members regularly supported one another when challenges arose and responsibilities were occasionally adjusted to ensure that progress continued smoothly. Collaborative tasks such as code reviews and system testing were shared among members to maintain quality and coherence across the project.

As the project deadline approached, the team reorganised into smaller subgroups responsible for different areas of development. This decision was motivated by the need to increase parallel work and reduce bottlenecks in critical tasks. Dividing responsibilities allowed specialised attention to areas such as testing, deployment and documentation, while still maintaining collaboration across the team. The allocation of responsibilities is summarised below:

\begin{table}[h]
\centering
\footnotesize
\caption{Subgroup Task Allocation Near Project Deadline}
\begin{tabularx}{\textwidth}{|p{4cm}|L|}
\hline
\textbf{Subgroup Task} & \textbf{Team Members} \\
\hline
Testing & Hafsa, Tejas, Benjamin \\
\hline
UI Fixes & Nitin, Ved \\
\hline
Deployment & Amey, Tejas \\
\hline
Report & Jasmeen, Khushi, Nitin \\
\hline
Refactoring & Amey, Nitin \\
\hline
\end{tabularx}
\label{tab:task-allocation}
\end{table}

This structure allowed the team to tackle multiple critical tasks simultaneously while maintaining collaboration and flexibility. By splitting into subgroups, each area of the project received focused attention, which helped ensure the project was completed on time and met the quality standards set by the team.

\section{Communication and Collaboration}
Communication played a critical role in the success of the project. The team held scheduled meetings twice a week (every Monday and Thursday) to review progress, discuss completed tasks and plan the next stage of development. Between meetings, informal communication was maintained via a WhatsApp group chat to share updates, ask questions and to resolve minor issues promptly.
In addition to meetings and group messaging, the team used a shared task board to coordinate development activities and maintain visibility of project progress. Code reviews were conducted regularly to maintain high standards of quality. Automated testing was implemented strategically to check functionality without affecting existing code.

\section{Challenges Encountered}
Several challenges arose during the project, which required careful planning and collaboration to overcome. One significant challenge was coordinating meetings, as team members had differing schedules. To address this, a hybrid approach was adopted: those who could attend in person participated in face-to-face meetings, while those who could not joined via WhatsApp calls. 

Another challenge was the use of React, as some team members were not familiar with the framework, requiring additional learning time and peer support. The team also encountered difficulties with overly ambitious ideas, which needed to be scaled down to ensure the Minimum Viable Product could be delivered on time.

Additional challenges included merging code and managing version control conflicts in Git, which occasionally slowed progress. Testing and integrating multiple components required more time than initially anticipated, particularly for complex features. Furthermore, the team had to navigate time pressure as deadlines approached, while maintaining quality standards. Clear communication and collaboration were critical in resolving any misunderstandings or minor conflicts that arose during the development process.

To overcome these challenges, the team refined the project scope, prioritised essential features and allocated tasks according to individual strengths. Less critical features were postponed or simplified, ensuring that the core system was completed successfully within the deadline. The hybrid meeting strategy, combined with collaborative support and task management, enabled the team to stay coordinated and productive despite these difficulties. A key factor that helped the team navigate these challenges was maintaining flexible communication channels.

\section{Lessons Learned}
Throughout this project, the team gained valuable experience in collaboration and adapting to unexpected challenges.

One of the key lessons was the importance of flexible communication; hybrid meetings, which combined in-person attendance with WhatsApp calls, proved essential for accommodating differing schedules and ensuring everyone remained involved.

The team also learned the value of clearly defining a Minimum Viable Product (MVP) and prioritising tasks, which helped keep the project manageable and ensured that critical features were delivered on time.

Working collaboratively on the codebase and coordinating development tasks helped the team improve both technical and problem-solving skills, particularly when learning new technologies such as React.

The project also reinforced important lessons about managing ambition, balancing the desire to implement additional features with the need to meet deadlines, and maintaining a positive and supportive team environment under time pressure. Overall, these experiences provided valuable insight into effective project management, communication, and teamwork, which will guide the team in future projects and professional settings.

Although the team maintained consistent communication and coordination, earlier adoption of more structured sprint planning and task estimation could have improved workload prediction. In several instances, tasks took longer than expected, which increased time pressure toward the project's final stages. More formal estimation techniques could have helped manage this more effectively and improve planning accuracy in future projects.